% GENERATED FILE. Edit canonical lessons, then run npm run generate:books.
% canonical-source-hash: fnv1a64:740765f404af9b63
% canonical-lessons: JA-C524-katei, JA-C524-omoiegaku, JA-C524-souzou-suru, JA-C524-suiri, JA-C524-jitsugen-suru

\chapter{Imagining}
\label{ch:ja-imagining}
\hlchaptermodality{524}

\begin{chapteropening}
\textbf{By the end of this chapter:} \emph{I can talk about what might be.}

Complete the last lesson of chapter 524: I can talk about what might be.
\end{chapteropening}

\section[\texorpdfstring{\ja{かてい} (katei)}{かてい (katei)}]{\ja{かてい} (katei) — a supposition; a hypothesis}

\label{lesson:JA-C524-katei}



\begin{quote}
Before the new one: say the Japanese for as it happens, then the Japanese for somehow (or other).
\end{quote}

\subsection*{You'll want to know: \ja{かてい}}
\textbf{\ja{かてい}} — \emph{katei} — \textquotedblleft{}a supposition; a hypothesis\textquotedblright{}.

\begin{quote}
\textbf{\ja{こたえは} \ja{かていに} \ja{よります。}} — \emph{kotae wa katei ni yorimasu} — \textquotedblleft{}The answer depends on the assumption.\textquotedblright{}
\end{quote}

\begin{cousinweb}
\textbf{\ja{かてい}} (\emph{katei}) is Sino-Japanese (kango): two Chinese characters read with their on'yomi, readings borrowed from Chinese. They are \emph{ka}, \textquotedblleft{}temporary\textquotedblright{} (Middle Chinese \emph{kaeX}) and \emph{tei}, \textquotedblleft{}determine\textquotedblright{} (\emph{dengH}) (KANJIDIC2; Middle Chinese in Baxter's transcription).
\end{cousinweb}

\begin{grammarlens}[title={What you've built}]
One more word for saying what happens if.
\end{grammarlens}

\subsection*{Your turn}
\begin{itemize}
\raggedright
  \item \emph{Say it:} \emph{katei}
  \item \emph{Say it:} \emph{katei}, once more
  \item \emph{Say it:} \emph{kotae wa katei ni yorimasu} — The answer depends on the assumption.
  \item \emph{Recall:} say the Japanese for as it happens, then the Japanese for somehow (or other), then say \emph{katei} again
\end{itemize}

\begin{tcolorbox}[breakable,colback=teal!4,colframe=teal!35!black,title={Before you move on}]
What does \ja{かてい} mean? (\textquotedblleft{}A supposition; a hypothesis\textquotedblright{}.) Say it once more.
\end{tcolorbox}
\section[\texorpdfstring{\ja{おもいえがく} (omoiegaku)}{おもいえがく (omoiegaku)}]{\ja{おもいえがく} (omoiegaku) — to picture (in one's mind)}

\label{lesson:JA-C524-omoiegaku}



\begin{quote}
Before the new one: say the Japanese for somehow (or other), then the Japanese for a supposition.
\end{quote}

\subsection*{You'll want to know: \ja{おもいえがく}}
\textbf{\ja{おもいえがく}} — \emph{omoiegaku} — \textquotedblleft{}to picture (in one's mind)\textquotedblright{}.

\begin{quote}
\textbf{\ja{もし} \ja{しょうらいを} \ja{おもいえがいたら、なにが} \ja{みえますか。}} — \emph{moshi shōrai o omoiegaitara, nani ga miemasu ka} — \textquotedblleft{}If you picture your future, what do you see?\textquotedblright{}
\end{quote}

Polite \textbf{\ja{おもいえがきます}}, \emph{omoiegakimasu}; a godan verb. Past (polite): \textbf{\ja{かれが} \ja{しょうらいを} \ja{おもいえがきました。}}, \emph{kare ga shōrai o omoiegakimashita}.

\begin{cousinweb}
\textbf{\ja{おもいえがく}} (\emph{omoiegaku}) is a native Japanese word (wago), inherited rather than borrowed from Chinese; its kanji was attached later as a kun'yomi, and its older origin is not settled (JMdict, KANJIDIC2).
\end{cousinweb}

\begin{grammarlens}[title={What you've built}]
One more word for saying what happens if.
\end{grammarlens}

\subsection*{Your turn}
\begin{itemize}
\raggedright
  \item \emph{Say it:} \emph{omoiegaku}
  \item \emph{Say it:} \emph{omoiegaku}, once more
  \item \emph{Say it:} \emph{moshi shōrai o omoiegaitara, nani ga miemasu ka} — If you picture your future, what do you see?
  \item \emph{Recall:} say the Japanese for somehow (or other), then the Japanese for a supposition, then say \emph{omoiegaku} again
\end{itemize}

\begin{tcolorbox}[breakable,colback=teal!4,colframe=teal!35!black,title={Before you move on}]
What does \ja{おもいえがく} mean? (\textquotedblleft{}To picture (in one's mind)\textquotedblright{}.) Say it once more.
\end{tcolorbox}
\section[\texorpdfstring{\ja{そうぞうする} (sōzō suru)}{そうぞうする (sōzō suru)}]{\ja{そうぞうする} (sōzō suru) — to imagine}

\label{lesson:JA-C524-souzou-suru}



\begin{quote}
Before the new one: say the Japanese for a supposition, then the Japanese for to picture (in one's mind).
\end{quote}

\subsection*{You'll want to know: \ja{そうぞうする}}
\textbf{\ja{そうぞうする}} — \emph{sōzō suru} — \textquotedblleft{}to imagine\textquotedblright{}.

\begin{quote}
\textbf{\ja{もし} \ja{こどもが} \ja{えほんを} \ja{よんだら、つきの} \ja{まちを} \ja{そうぞうするでしょう。}} — \emph{moshi kodomo ga ehon o yondara, tsuki no machi o sōzō suru deshō} — \textquotedblleft{}If children read the picture book, they would imagine a town on the moon.\textquotedblright{}
\end{quote}

Polite \textbf{\ja{そうぞうします}}, \emph{sōzōshimasu}; a \ja{する} verb. Past (polite): \textbf{\ja{こどもが} \ja{つきの} \ja{まちを} \ja{そうぞうしました。}}, \emph{kodomo ga tsuki no machi o sōzōshimashita}.

\begin{cousinweb}
\textbf{\ja{そうぞうする}} (\emph{sōzō suru}) is the Sino-Japanese (kango) noun \emph{sōzō} with the native verb \emph{suru}, \textquotedblleft{}to do\textquotedblright{}: two Chinese characters read with their on'yomi, readings borrowed from Chinese. They are \emph{sō}, \textquotedblleft{}concept\textquotedblright{} (Middle Chinese \emph{sjangX}) and \emph{zō}, \textquotedblleft{}statue\textquotedblright{} (\emph{zjangX}) (KANJIDIC2; Middle Chinese in Baxter's transcription).
\end{cousinweb}

\begin{grammarlens}[title={What you've built}]
One more word for saying what happens if.
\end{grammarlens}

\subsection*{Your turn}
\begin{itemize}
\raggedright
  \item \emph{Say it:} \emph{sōzō suru}
  \item \emph{Say it:} \emph{sōzō suru}, once more
  \item \emph{Say it:} \emph{moshi kodomo ga ehon o yondara, tsuki no machi o sōzō suru deshō} — If children read the picture book, they would imagine a town on the moon.
  \item \emph{Recall:} say the Japanese for a supposition, then the Japanese for to picture (in one's mind), then say \emph{sōzō suru} again
\end{itemize}

\begin{tcolorbox}[breakable,colback=teal!4,colframe=teal!35!black,title={Before you move on}]
What does \ja{そうぞうする} mean? (\textquotedblleft{}To imagine\textquotedblright{}.) Say it once more.
\end{tcolorbox}
\section[\texorpdfstring{\ja{すいり} (suiri)}{すいり (suiri)}]{\ja{すいり} (suiri) — reasoning; deduction}

\label{lesson:JA-C524-suiri}



\begin{quote}
Before the new one: say the Japanese for to picture (in one's mind), then the Japanese for to imagine.
\end{quote}

\subsection*{You'll want to know: \ja{すいり}}
\textbf{\ja{すいり}} — \emph{suiri} — \textquotedblleft{}reasoning; deduction\textquotedblright{}.

\begin{quote}
\textbf{\ja{もし} \ja{かれの} \ja{すいりが} \ja{あたったら、すごいでしょう。}} — \emph{moshi kare no suiri ga atattara, sugoi deshō} — \textquotedblleft{}If his reasoning turns out right, that will be amazing.\textquotedblright{}
\end{quote}

\begin{cousinweb}
\textbf{\ja{すいり}} (\emph{suiri}) is Sino-Japanese (kango): two Chinese characters read with their on'yomi, readings borrowed from Chinese. They are \emph{sui}, \textquotedblleft{}conjecture\textquotedblright{} (Middle Chinese \emph{tsyhwij}) and \emph{ri}, \textquotedblleft{}logic\textquotedblright{} (\emph{liX}) (KANJIDIC2; Middle Chinese in Baxter's transcription).
\end{cousinweb}

\begin{grammarlens}[title={What you've built}]
One more word for saying what happens if.
\end{grammarlens}

\subsection*{Your turn}
\begin{itemize}
\raggedright
  \item \emph{Say it:} \emph{suiri}
  \item \emph{Say it:} \emph{suiri}, once more
  \item \emph{Say it:} \emph{moshi kare no suiri ga atattara, sugoi deshō} — If his reasoning turns out right, that will be amazing.
  \item \emph{Recall:} say the Japanese for to picture (in one's mind), then the Japanese for to imagine, then say \emph{suiri} again
\end{itemize}

\begin{tcolorbox}[breakable,colback=teal!4,colframe=teal!35!black,title={Before you move on}]
What does \ja{すいり} mean? (\textquotedblleft{}Reasoning; deduction\textquotedblright{}.) Say it once more.
\end{tcolorbox}
\section[\texorpdfstring{\ja{じつげんする} (jitsugen suru)}{じつげんする (jitsugen suru)}]{\ja{じつげんする} (jitsugen suru) — to be realized; to come true}

\label{lesson:JA-C524-jitsugen-suru}



\begin{quote}
Before the new one: say the Japanese for to imagine, then the Japanese for reasoning.
\end{quote}

\subsection*{You'll want to know: \ja{じつげんする}}
\textbf{\ja{じつげんする}} — \emph{jitsugen suru} — \textquotedblleft{}to be realized; to come true\textquotedblright{}.

\begin{quote}
\textbf{\ja{もし} \ja{ゆめが} \ja{じつげんしたら、あねは} \ja{よろこぶでしょう。}} — \emph{moshi yume ga jitsugen shitara, ane wa yorokobu deshō} — \textquotedblleft{}If her dream comes true, my elder sister will be overjoyed.\textquotedblright{}
\end{quote}

Polite \textbf{\ja{じつげんします}}, \emph{jitsugenshimasu}; a \ja{する} verb. Past (polite): \textbf{\ja{かれの} \ja{ゆめが} \ja{じつげんしました。}}, \emph{kare no yume ga jitsugenshimashita}.

\begin{cousinweb}
\textbf{\ja{じつげんする}} (\emph{jitsugen suru}) is the Sino-Japanese (kango) noun \emph{jitsugen} with the native verb \emph{suru}, \textquotedblleft{}to do\textquotedblright{}: two Chinese characters read with their on'yomi, readings borrowed from Chinese. They are \emph{jitsu}, \textquotedblleft{}reality\textquotedblright{} (Middle Chinese \emph{zyit}) and \emph{gen}, \textquotedblleft{}present\textquotedblright{} (\emph{henH}) (KANJIDIC2; Middle Chinese in Baxter's transcription).
\end{cousinweb}

\begin{grammarlens}[title={What you've built}]
One more word for saying what happens if.
\end{grammarlens}

\subsection*{Your turn}
\begin{itemize}
\raggedright
  \item \emph{Say it:} \emph{jitsugen suru}
  \item \emph{Say it:} \emph{jitsugen suru}, once more
  \item \emph{Say it:} \emph{moshi yume ga jitsugen shitara, ane wa yorokobu deshō} — If her dream comes true, my elder sister will be overjoyed.
  \item \emph{Recall:} say the Japanese for to imagine, then the Japanese for reasoning, then say \emph{jitsugen suru} again
\end{itemize}

\begin{tcolorbox}[breakable,colback=teal!4,colframe=teal!35!black,title={Before you move on}]
What does \ja{じつげんする} mean? (\textquotedblleft{}To be realized; to come true\textquotedblright{}.) Say it once more.
\end{tcolorbox}
